\section{Problem Description}\label{sec:problem}
%==============================
\subsection{Workspace and Robots}\label{subsec:ws}
The workspace is a bounded planar region $\mathcal{W}\subset\mathbb{R}^2$
containing immovable structures~$\mathcal{O}\subset\mathcal{W}$ and $M\geq0$
movable rigid polygons
$\boldsymbol{\Omega}\triangleq\{\Omega_1,\cdots,\Omega_M\}\subset\mathcal{W}$.
Each movable obstacle~$\Omega_m$ has mass~$\mathsf{M}_m$,
inertia~$\mathsf{I}_m$, frictional parameters and state
$\mathbf{s}_m(t)\triangleq(\mathbf{x}_m(t),\psi_m(t))$, with
$\mathbf{x}_m(t)\in\mathbb{R}^2$ the centroid position and
$\psi_m(t)\in\mathbb{R}$ the orientation; $\Omega_m(t)$ denotes the region it
occupies at time~$t$. A team of~$N$ robots,
$\mathcal{R}\triangleq\{R_1,\cdots,R_N\}$, operates as a cooperative unit. Each
robot~$R_i$ is a rigid body with state
$\mathbf{s}_{R_i}(t)\triangleq(\mathbf{x}_{R_i}(t),\psi_{R_i}(t))$ and
footprint~$R_i(t)$. The instantaneous free space is given by:
\begin{equation}\label{eq:freespace}
\widehat{\mathcal{W}}(t)\triangleq\mathcal{W}\setminus\Big(
\mathcal{O}\cup\{\Omega_m(t)\}_{m=1}^{M}\cup\{R_i(t)\}_{i=1}^{N}
\Big),
\end{equation}
where~$\widehat{\mathcal{W}}(t)$ excludes regions occupied by immovable
obstacles, movable obstacles, and robots.

%==============================
\subsection{External Vehicle and Clearance Goal}\label{subsec:vehicle}
An external vehicle~$\texttt{V}$ whose footprint is enclosed by a disk of
radius~$W/2>0$ must navigate within the workspace from a start position
$\mathbf{x}_\texttt{V}^{\texttt{S}}\in\mathcal{W}$ to a goal position
$\mathbf{x}_\texttt{V}^{\texttt{G}}\in\mathcal{W}$. Since movable obstacles may
obstruct the way, the clearance goal is to rearrange them so that the workspace
contains a passage from~$\mathbf{x}_\texttt{V}^{\texttt{S}}$ to
$\mathbf{x}_\texttt{V}^{\texttt{G}}$ that can accommodate this enclosing disk,
i.e., a passage of width at least~$W$. At time~$t$, this is expressed by the
admissible reference-center space
$\mathcal{F}_W(t)\triangleq
\{p\in\widehat{\mathcal{W}}(t)\mid
\mathbb{B}_{W/2}(p)\subset\widehat{\mathcal{W}}(t)\}$, where
$\mathbb{B}_{W/2}(p)$ is the disk of radius~$W/2$ centered at~$p$. The
$W$-clearance condition is then given by:
\begin{equation}\label{eq:wclear}
\exists\,\mathcal{P}^W_\texttt{V}\subset\mathcal{F}_W(t):\
\mathbf{x}_\texttt{V}^{\texttt{S}}\leadsto\mathbf{x}_\texttt{V}^{\texttt{G}},
\end{equation}
which certifies a collision-free centerline path for any vehicle footprint
contained in the enclosing disk.
% Since movable obstacles may obstruct the way,
% a direct passage is not always feasible.
% At time~$t$, the admissible center positions of~$\texttt{V}$ are
% \begin{equation}
% \mathcal{F}_W(t)\triangleq
% \{p\in\widehat{\mathcal{W}}(t)\mid \mathbb{B}_{W/2}(p)\subset\widehat{\mathcal{W}}(t)\}.
% \end{equation}
% where $\mathbb{B}_{W/2}$ is a disk of radius~$W/2$.
% The clearance goal is satisfied if there exists a continuous centerline path
% inside~$\mathcal{F}_W(t)$:
% \begin{equation}\label{eq}
% \exists\,\mathcal{P}^W_\texttt{V}\subset\mathcal{F}_W(t):
% \mathbf{x}_\texttt{V}^{\texttt{S}}\leadsto\mathbf{x}_\texttt{V}^{\texttt{G}}.
% \end{equation}
% this therefore certifies a passage with width at
% least~$W$.


%==============================


\subsection{Collaborative Pushing Modes}\label{ss:interaction_mode}
To enlarge blocked gaps, the robots may actively
reconfigure~$\boldsymbol{\Omega}$ by pushing movable obstacles. For obstacle
$\Omega_m\in\boldsymbol{\Omega}$, a pushing mode is denoted by
$\boldsymbol{\xi}_m\triangleq(\mathbf{c}_m,\mathbf{u}_m)$, where
$\mathbf{c}_m=(c_1,\cdots,c_N)$ is an ordered tuple of contact points on
$\partial\Omega_m$, and the ordering specifies the robot-to-contact assignment.
The vector~$\mathbf{u}_m\in\mathbb{R}^{2N}$ stacks the corresponding planar
contact forces, expressed in the body frame of~$\Omega_m$.
For a contact tuple~$\mathbf{c}_m$, the set of admissible forces
${U}_m(\mathbf{c}_m)\subset\mathbb{R}^{2N}$ collects contact forces
satisfying the robot pushing limits, robot--object friction, and unilateral
contact constraints. The body wrench induced by a mode is defined as:
\begin{equation}
\mathbf{q}_m(\boldsymbol{\xi}_m)
\triangleq
G_m(\mathbf{c}_m)\mathbf{u}_m,
\label{eq:contact_wrench_map}
\end{equation}
where~$G_m(\mathbf{c}_m)\in\mathbb{R}^{3\times 2N}$ is the contact-to-wrench mapping;
and $\mathbf{q}_m(\boldsymbol{\xi}_m)\in\mathbb{R}^3$.
The admissible mode set is state dependent as denoted by~$\Xi_m(\mathbf{s})$.
It contains modes whose contacts are reachable from the
current robot configurations, pre-contact configurations are feasible in
the current clutter, and forces satisfy
$\mathbf{u}_m\in{U}_m(\mathbf{c}_m)$. Hence, applying a feasible
mode~$\boldsymbol{\xi}_m\in\Xi_m(\mathbf{s}(t))$ yields the physical
transition:
\begin{equation}\label{eq:transition}
  \mathbf{s}(t^+)\triangleq
  \Phi\big(\mathbf{s}(t),\,m,\,\boldsymbol{\xi}_m\big),
\end{equation}
where~$\mathbf{s}(t)$ stacks all robot and obstacle states.
\change{
The force vector~$\mathbf{u}_m$ serves as a feasibility variable to certify
wrench generation under the modeled contact and actuation constraints rather
than a directly commanded control input.
}


%==============================
\subsection{Problem Statement}\label{subsec:objective}

The goal is to compute a \emph{hybrid schedule} that reconfigures movable
obstacles to admit a $W$-feasible path for the external vehicle.
The schedule is defined as:
\begin{equation}
    \pi\triangleq
    \left\{
    \left(
    \Omega_{m_k},
    \boldsymbol{\xi}_k,
    \Delta t_k
    \right)
    \right\}_{k=1}^{K},
\end{equation}
where~$\Omega_{m_k}\in\boldsymbol{\Omega}$ is the manipulated obstacle,
$\boldsymbol{\xi}_k\in\Xi_{m_k}$ is its pushing mode, and
$\Delta t_k>0$ is the execution duration. The search is guided by
\begin{equation}\label{eq:problem}
    \mathcal{J}_{\mathrm{clr}}(\pi)
    \triangleq
    T+
    \alpha
    \sum_{k=1}^{K}
    J\left(
    m_k,\boldsymbol{\xi}_k,\mathbf{s}(t_k)
    \right),
\end{equation}
where~$T\triangleq\sum_{k=1}^{K}\Delta t_k$ denotes the robotic clearance
duration excluding subsequent vehicle traversal, and~$J(\cdot)$ is the
simulation-estimated manipulation effort. Vehicle traversal is considered
only through the terminal $W$-clearance condition~\eqref{eq:wclear}.
The constraints~\eqref{eq:freespace}--\eqref{eq:transition} enforce valid
pushing transitions. \change{Equation~\eqref{eq:problem} ranks feasible
schedules; PIHS uses finite sampling and does not guarantee global optimality of
the continuous hybrid objective.}

% ==============================
\begin{remark}\label{remark:uniqueness}
  \change{The objective and constraints above couple obstacle selection, pushing-mode generation, and switching in one hybrid planning problem, yielding a combinatorial-continuous search space far richer than classical NAMO formulations with simple objects or hand-crafted contact modes~\cite{goyal1989limit,chen2015occlusion,wang2006multi}.}

  \hfill~$\blacksquare$
\end{remark}
% ==============================
